\documentclass[acmsmall,screen,nonacm]{acmart}

\usepackage{macros}
\usepackage{mathpartir}
\usepackage{cleveref}

\begin{document}

\title{Operational Semantics of the \textsf{alifib} Interpreter}
\author{} \date{}
\maketitle

\begin{abstract}
We give a big-step operational semantics for the \textsf{alifib}
interpreter. The language manipulates \emph{directed complexes}---oriented
graded posets labelled by cell generators and equipped with partial maps
between them---for the purpose of formalising constructions in higher
category theory. We describe the abstract syntax, the interpreter state,
and give evaluation rules for every language construct. We draw attention
to several subtle design decisions and their rationale.
\end{abstract}

\tableofcontents

% ===========================================================================
\section{Introduction}
\label{sec:intro}
% ===========================================================================

An \textsf{alifib} program declares \emph{types} (each of which carries a
directed complex of generators), \emph{cells} (generators at various
dimensions with specified source and target boundaries), \emph{diagrams}
(labelled pasting diagrams built by composition of cells), and
\emph{partial maps} (structure-preserving maps between complexes).
Intuitively, a directed complex is an oriented graded poset---a
collection of cells at each dimension, with face maps connecting each cell
to the lower-dimensional cells on its boundary---that serves as the
combinatorial data underlying a pasting diagram in a higher category.
The interpreter maintains a mutable global state and threads it through
the evaluation of each instruction.

This document provides a self-contained formal account of the
interpreter's behaviour, abstracting away Rust-level implementation
details (error reporting, span tracking, \texttt{Arc}-based sharing)
while preserving every semantically meaningful step.

\paragraph{Conventions.}
We write $\State$ for the global state, $\Complex$ for a complex (the
local namespace of generators, diagrams, and maps associated with a type
or module), and $n$ for names. We use $\Downarrow$ for big-step
evaluation. Metavariables are colour-coded: \Type{types} in blue,
\Cell{cells} in brown, \Map{maps} in teal, and \Diag{diagrams} in
purple. Keywords appear in $\kw{sans\text{-}serif}$.
All nonterminals in the grammar are set in \textsf{sans-serif}.

% ===========================================================================
\section{Abstract Syntax}
\label{sec:syntax}
% ===========================================================================

We present the abstract syntax in layers. We simplify the concrete syntax
wherever it would introduce unnecessary detail: for instance, we omit span
annotations and flatten optional fields. The concrete syntax of
\textsf{alifib} can be recovered by consulting the parser.

% ---------------------------------------------------------------------------
\subsection{Programs, Blocks, and Instructions}
\label{sec:syntax:prog}
% ---------------------------------------------------------------------------

A program is a sequence of blocks. There are two kinds: \emph{type blocks}
(introduced by $\kwtype$) for declaring new types, diagrams, maps, and
module inclusions at the top level; and \emph{local blocks} (introduced by
$\mathsf{@}\Cpx$) for working within an existing type's complex.

\begin{figure}[t]
\fbox{\text{programs and blocks}}
\[
  \begin{array}{rcll}
    \Prog  & \Coloneqq & \overline{\Block}         \\[3pt]
    \Block & \Coloneqq & \kwtype\;\overline{\TyInst}
             & \text{type block} \\
           &           & \mathsf{@}\Cpx\;\overline{\LocInst}
             & \text{local block}
  \end{array}
\]

\fbox{\text{type-level instructions}}
\[
  \begin{array}{rcll}
    \TyInst & \Coloneqq & n \;\decl\; \Cpx
              & \text{type declaration} \\
            &           & \kwlet\;n = \Diagram
              & \text{diagram naming} \\
            &           & \kwlet\;n \mathbin{::} \Addr = \PMapDef
              & \text{map definition} \\
            &           & \kwlet\;\kwtotal\;n \mathbin{::} \Addr = \PMapDef
              & \text{total map definition} \\
            &           & \kwinclude\;n\;[\kwas\;n']
              & \text{module inclusion}
  \end{array}
\]

\fbox{\text{complex-body instructions}}
\[
  \begin{array}{rcll}
    \CInst & \Coloneqq & n                   & \text{0-cell} \\
           &           & n : \Diagram \to \Diagram & \text{higher cell} \\
           &           & \kwlet\;n = \Diagram & \text{diagram naming} \\
           &           & \kwlet\;[\kwtotal]\;n \mathbin{::} \Addr = \PMapDef
             & \text{map definition} \\
           &           & \kwinclude\;\Addr\;[\kwas\;n]
             & \text{type inclusion} \\
           &           & \kwattach\;n \mathbin{::} \Addr\;[\kwalong\;\PMapDef]
             & \text{type attachment}
  \end{array}
\]

\fbox{\text{local instructions}}
\[
  \begin{array}{rcll}
    \LocInst & \Coloneqq & \kwlet\;n = \Diagram
               & \text{diagram naming} \\
             &           & \kwlet\;[\kwtotal]\;n \mathbin{::} \Addr = \PMapDef
               & \text{map definition} \\
             &           & \kwassert\;\Diagram = \Diagram
               & \text{equality assertion}
  \end{array}
\]
\caption{Programs and instructions.}
\label{fig:syntax:prog}
\end{figure}

\Cref{fig:syntax:prog} gives the full grammar of programs and
instructions. All forms that define partial maps use the $\kwlet$ keyword
in the concrete syntax; the distinction between diagram bindings and map
definitions is syntactic (map definitions include a $\mathbin{::} \Addr$
annotation). The keyword $\kwtotal$ marks a map definition as total,
meaning the interpreter checks that every generator in the source complex
has an image.

The type declaration form $n \;\decl\; \Cpx$ introduces a new 0-dimensional
type with a complex body. The symbol $\decl$ is written \texttt{<<=} in
the concrete syntax. The interpreter rejects type declarations with
non-trivial boundaries: all types are 0-cells, and higher cells are
declared inside complex bodies.

% ---------------------------------------------------------------------------
\subsection{Complexes and Addresses}
\label{sec:syntax:cpx}
% ---------------------------------------------------------------------------

A complex reference identifies which type's complex to operate in:
\[
  \begin{array}{rcll}
    \Cpx  & \Coloneqq & \Addr
            & \text{resolved address} \\
          &           & \Addr\;\{ \overline{\CInst} \}
            & \text{address with inline body} \\[3pt]
    \Addr & \Coloneqq & \varepsilon \mid n_1.n_2.\cdots.n_k
  \end{array}
\]
The empty address $\varepsilon$ refers to the current module's anonymous root
type. A non-empty address is resolved by navigating module-domain maps
for all segments except the last, and then looking up the final segment
as a cell diagram (\cref{sec:complex:addr}).

% ---------------------------------------------------------------------------
\subsection{Diagrams}
\label{sec:syntax:diag}
% ---------------------------------------------------------------------------

\begin{figure}[t]
\fbox{\text{diagram expressions}}
\[
  \begin{array}{rcll}
    \Diagram & \Coloneqq & \DExpr_1 \;\cdots\; \DExpr_k
               & \text{implicit pasting} \\
             &           & \Diagram \paste{k} \DExpr_1 \;\cdots\; \DExpr_j
               & \text{explicit pasting at dimension $k$} \\[3pt]
    \DExpr   & \Coloneqq & \DComp
               & \text{component} \\
             &           & \DExpr.\DComp
               & \text{dot access} \\[3pt]
    \DComp   & \Coloneqq & n
               & \text{name lookup} \\
             &           & (\Diagram)
               & \text{parenthesised diagram} \\
             &           & (\PMapG)
               & \text{parenthesised map} \\
             &           & (\,\kwmap\;\PMapDef \mathbin{::} \Cpx\,)
               & \text{anonymous map} \\
             &           & \kwin \mid \kwout
               & \text{boundary accessors} \\
             &           & \kwhole
               & \text{hole}
  \end{array}
\]

\fbox{\text{partial map expressions}}
\[
  \begin{array}{rcll}
    \PMapDef  & \Coloneqq & \PMapG
                & \text{plain map} \\
              &           & \PMapExt
                & \text{extension map} \\[3pt]
    \PMapG    & \Coloneqq & \PMapBasic
                & \text{basic map} \\
              &           & \PMapBasic.\PMapG
                & \text{composition via dot} \\[3pt]
    \PMapBasic & \Coloneqq & n
                & \text{named map} \\
              &           & (\,\kwmap\;\PMapDef \mathbin{::} \Cpx\,)
                & \text{anonymous map} \\
              &           & (\PMapG)
                & \text{parenthesised} \\[3pt]
    \PMapExt  & \Coloneqq & [\,\overline{\Clause}\,]
                & \text{extension from empty} \\
              &           & \PMapG\;[\,\overline{\Clause}\,]
                & \text{extension of prefix} \\[3pt]
    \Clause   & \Coloneqq & \Diagram \Rightarrow \Diagram
  \end{array}
\]
\caption{Diagram and partial map expressions.}
\label{fig:syntax:expr}
\end{figure}

\Cref{fig:syntax:expr} gives the grammar of expressions. There are two
kinds of pasting. \emph{Implicit pasting} (juxtaposition) pastes a
sequence of diagram expressions left to right, automatically computing the
paste dimension from the dimensions of the operands. \emph{Explicit
pasting} $\Diagram \paste{k} \overline{\DExpr}$ pastes at a specified
dimension $k$.

\paragraph{Name ambiguity.}
When a name $n$ is looked up in a complex, the interpreter checks the
diagram namespace first, then the map namespace. A name found as a diagram
produces a diagram term; a name found as a map produces a map term
carrying both the map and a reference to its source complex. This
overloading is intentional: the dot operator~(.) behaves differently
depending on whether its left operand is a diagram or a map
(\cref{sec:diag:dot}).

\paragraph{Holes.}
The hole symbol $\kwhole$ (written \texttt{?} in concrete syntax) is not
a first-class value. When the interpreter encounters a hole, it records
diagnostic information (including boundary context inferred from adjacent
expressions) and aborts diagram evaluation. Holes exist to support
interactive development.

% ===========================================================================
\section{Interpreter State}
\label{sec:state}
% ===========================================================================

The interpreter maintains a global state $\State$ consisting of three
finite maps:
\[
  \State = \langle \cells, \types, \modules \rangle
\]

\begin{itemize}
  \item $\cells : \GId \rightharpoonup (\mathbb{N} \times \CData)$ maps
    global identifiers to their dimension and cell data.
  \item $\types : \GId \rightharpoonup (\CData \times \Complex)$ maps
    global identifiers of types to their cell data and associated complex.
  \item $\modules : \textsf{ModId} \rightharpoonup \Complex$ maps canonical
    file paths to module-level complexes.
\end{itemize}

% ---------------------------------------------------------------------------
\subsection{Global Identifiers and Tags}
\label{sec:state:gid}
% ---------------------------------------------------------------------------

The state uses \emph{global identifiers} $\GId \in \mathbb{N}$ to name
cells and types across the entire program. Identifiers are allocated by a
monotonic counter: $\fresh$ returns the next unused natural number. Every
globally-registered cell or type receives a unique $\GId$ at the point of
declaration.

The purpose of global identifiers is to provide a stable, unique identity
for each cell that is invariant under renaming. Two cells in different
complexes may bear the same human-readable name (e.g.\ \texttt{pt}) but
refer to entirely different generators; conversely, the \emph{same} cell
may appear in multiple complexes under different names (e.g.\ after an
$\kwinclude$). Global identifiers resolve this: they are the canonical
identity of a cell. In the Rust implementation, $\GId$ is a wrapper around
\texttt{usize} allocated via a global atomic counter.

Every cell in a diagram is labelled not with a name but with a \emph{tag}:
\[
  \textsf{Tag} \Coloneqq \TagG{\GId} \mid \TagL{n}
\]
Global tags $\TagG{\GId}$ reference cells registered in $\State.\cells$
or $\State.\types$. Local tags $\TagL{n}$ reference cells in a complex's
local cell store; they are used only inside complex bodies being
constructed and must not appear in persisted definitions. The reason for
having two kinds of tags is that local cells are temporary scaffolding
during complex construction: they exist in the working copy of a complex
but are not registered in the global state.

% ---------------------------------------------------------------------------
\subsection{Cell Data}
\label{sec:state:celldata}
% ---------------------------------------------------------------------------

Every cell has associated \emph{cell data} describing its boundary:
\[
  \CData \Coloneqq \CZero \mid \CBdy{d^{-}}{d^{+}}
\]
The symbol $\CZero$ denotes a 0-cell (no boundaries). The form
$\CBdy{d^{-}}{d^{+}}$ describes a cell of dimension $n > 0$ with
input boundary $d^{-}$ and output boundary $d^{+}$, both diagrams of
dimension $n - 1$.

% ---------------------------------------------------------------------------
\subsection{Complexes}
\label{sec:state:complex}
% ---------------------------------------------------------------------------

A \emph{complex} $\Complex$ is the namespace associated with a type or
module. It bundles four components:
\[
  \Complex = \langle \mathit{gens}, \mathit{diags}, \mathit{maps},
    \mathit{locals} \rangle
\]
\begin{itemize}
  \item $\mathit{gens} : n \rightharpoonup
    (\textsf{Tag} \times \mathbb{N} \times \Diag{D})$
    maps generator names to their tag, dimension, and
    \emph{classifier diagram} (the diagram consisting of that single cell).
  \item $\mathit{diags} : n \rightharpoonup \Diag{D}$ maps diagram names
    to diagram values.
  \item $\mathit{maps} : n \rightharpoonup
    (\textsf{MapDom} \times \Map{F})$ maps partial map names to a domain
    annotation and a map value.
    The domain is either $\mathsf{Ty}(\GId)$ or $\mathsf{Mod}(\textsf{ModId})$.
  \item $\mathit{locals} : n \rightharpoonup (\mathbb{N} \times \CData)$
    stores locally-scoped cells (populated only during complex body
    interpretation).
\end{itemize}
A complex also maintains a set of all names in use, to detect collisions
eagerly. When a generator is \emph{declared} in a complex body, it is
entered into both $\mathit{gens}$ and $\mathit{diags}$: the former
records structural metadata while the latter enables name resolution
during diagram evaluation. Generators imported via $\kwinclude$ or
$\kwattach$ are entered only into $\mathit{gens}$; they are accessed
through dot notation on the inclusion map rather than by direct name
lookup in $\mathit{diags}$.

% ---------------------------------------------------------------------------
\subsection{Terms}
\label{sec:state:terms}
% ---------------------------------------------------------------------------

Diagram expressions evaluate to \emph{terms}:
\[
  \TermD \Coloneqq \DTerm{\Diag{D}} \mid
  \MTerm{\Map{F}, \Complex_{\mathit{src}}}
\]
A diagram term $\DTerm{\Diag{D}}$ carries a concrete diagram. A map
term $\MTerm{\Map{F}, \Complex_{\mathit{src}}}$ carries a partial map
$\Map{F}$ together with a reference to its \emph{source complex}. The
source complex is essential: it determines the namespace in which further
dot-access into the map is resolved.

The reason the interpreter tracks terms rather than just diagrams is that
partial maps and diagrams share the same expression language. When the user
writes \texttt{F.x}, the interpreter must know that \texttt{F} is a map
in order to resolve \texttt{x} in the source complex of \texttt{F} (not
in the ambient complex). For example, if \texttt{F} is a map from a type
\texttt{Group} to the current type, then \texttt{F.mul} looks up
\texttt{mul} in \texttt{Group}'s complex and returns its image under
\texttt{F}. The term type makes this dispatch explicit.

% ---------------------------------------------------------------------------
\subsection{State Threading}
\label{sec:state:threading}
% ---------------------------------------------------------------------------

The state $\State$ is threaded through interpretation. The judgment form is:
\[
  \IntJ{\State}{\Complex}{\mathit{instr}}{\State'}
\]
reading: ``under state $\State$ and location $\Complex$, interpreting
instruction $\mathit{instr}$ yields a new state $\State'$.'' Instructions
within a block are processed sequentially: each receives the state produced
by the previous. Local blocks stop on first error; type blocks and complex
bodies process all instructions regardless of errors, accumulating
diagnostics.

In the Rust implementation, $\State$ is wrapped in \texttt{Arc<State>} and
mutated via \texttt{Arc::make\_mut}. This copy-on-write discipline is not
observable in the formal rules but explains why the implementation freely
clones contexts before potentially-failing operations.

% ===========================================================================
\section{Program Interpretation}
\label{sec:prog}
% ===========================================================================

\subsection{Module Initialisation}
\label{sec:prog:init}

When a module $m$ is first interpreted, its complex is initialised with an
anonymous root type:

\begin{mathpar}
  \inferrule*[lab=\textsc{Prog-Skip}]
    {\findmod{m} \neq \bot}
    {\IntJ{\State}{\cdot}{\Prog(m)}{\State}}
  \and
  \inferrule*[lab=\textsc{Prog-Init}]
    {\findmod{m} = \bot \\
     \GId_r = \fresh \\
     D_r = \CellOf{\TagG{\GId_r}}{\CZero} \\\\
     \Complex_0 = [\varepsilon \mapsto D_r] \\
     \State_1 = \State[\types(\GId_r) \mapsto (\CZero, \emptyset),\;
                        \modules(m) \mapsto \Complex_0] \\\\
     \IntJ{\State_1}{\Complex_0}{\overline{\Block}}{\State'}}
    {\IntJ{\State}{\cdot}{\Prog(m)}{\State'}}
\end{mathpar}

\textsc{Prog-Skip} ensures idempotent module loading: if $m$ is already
loaded, its interpretation is skipped entirely. This prevents
double-interpretation when two modules include the same dependency, but
it also means that the order of includes can affect what state a module
sees.

\textsc{Prog-Init} allocates a fresh $\GId_r$ for an anonymous root type,
creates its cell diagram $D_r$, and registers both the type (with an
empty complex) and the module (with a complex containing $D_r$ under the
empty name $\varepsilon$). The root type and the module complex are
\emph{distinct} objects: the module complex holds the root generator
\emph{reference}, while the root type's complex will hold generators
declared by the first $\kwtype$ block that has no explicit address.

\subsection{Block Dispatch}
\label{sec:prog:block}

\begin{mathpar}
  \inferrule*[lab=\textsc{Blk-Type}]
    {\IntJ{\State}{\findmod{m}}{\overline{\TyInst}}{\State'}}
    {\IntJ{\State}{\cdot}{\kwtype\;\overline{\TyInst}}{\State'}}
  \and
  \inferrule*[lab=\textsc{Blk-Local}]
    {\ResJ{\State}{\Cpx}{(\GId_r, \Complex)} \\
     \IntJ{\State}{\Complex}{\overline{\LocInst}}{\State'}}
    {\IntJ{\State}{\cdot}{\mathsf{@}\Cpx\;\overline{\LocInst}}{\State'}}
\end{mathpar}

A type block interprets its instructions within the module complex. A local
block first resolves its complex reference to a namespace
$(\GId_r, \Complex)$ and then interprets local instructions within that
complex.

% ===========================================================================
\section{Type-Level Instructions}
\label{sec:type}
% ===========================================================================

\subsection{Type Declaration}
\label{sec:type:gen}

\begin{mathpar}
  \inferrule*[lab=\textsc{TyDecl}]
    {n \notin \mathit{used}(\findmod{m}) \\
     \ResJ{\State}{\Cpx}{(\_, \Complex_{\mathit{def}})} \\
     \GId_T = \fresh \\
     D_T = \CellOf{\TagG{\GId_T}}{\CZero} \\\\
     F_{\mathit{id}} = \idmap{\Complex_{\mathit{def}}} \\
     \Complex'_{\mathit{def}} = \Complex_{\mathit{def}}
       [\mathit{maps}(n) \mapsto (\mathsf{Ty}(\GId_T), F_{\mathit{id}})] \\\\
     \State' = \State[\types(\GId_T) \mapsto
       (\CZero, \Complex'_{\mathit{def}}),\;
       \modules(m)(n) \mapsto D_T]}
    {\IntJ{\State}{\findmod{m}}{n \;\decl\; \Cpx}{\State'}}
\end{mathpar}

A type declaration resolves the complex body $\Cpx$ (interpreting any cell
generators, includes, and attaches inside), allocates a fresh $\GId_T$,
constructs the identity map $F_{\mathit{id}}$ on the definition complex,
and adds this identity map to the type's own complex under its own name.
The classifier diagram $D_T$ is added to both the generator and diagram
namespaces of the module complex.

\paragraph{The self-referencing identity map.}
The identity map is added to the type's own complex under the type's name.
This is how a type ``refers to itself'': inside the complex of a type
$T$, the name $T$ resolves to the identity inclusion map. This is
essential for constructions like $\kwattach\;S \mathbin{::} \textrm{Semigroup}\;\kwalong\;[\textrm{Set} \Rightarrow \textrm{Set}]$, where the right-hand \textrm{Set}
refers to the current type's own identity.

% ---------------------------------------------------------------------------
\subsection{Diagram Naming}
\label{sec:type:let}
% ---------------------------------------------------------------------------

\begin{mathpar}
  \inferrule*[lab=\textsc{TyLet}]
    {\EvalD{\State}{\findmod{m}}{\Diagram}{D}}
    {\IntJ{\State}{\findmod{m}}{\kwlet\;n = \Diagram}%
          {\State[\modules(m).\mathit{diags}(n) \mapsto D]}}
\end{mathpar}

The diagram expression is evaluated and bound in the module complex.
Unlike complex-body and local-block let-bindings, the Rust
implementation does \emph{not} check name uniqueness at the type level:
a duplicate name silently overwrites the previous binding.

% ---------------------------------------------------------------------------
\subsection{Map Definition}
\label{sec:type:def}
% ---------------------------------------------------------------------------

\begin{mathpar}
  \inferrule*[lab=\textsc{TyMap}]
    {\AddrJ{\State}{\Addr}{\GId_S} \\
     \Complex_S = \findtype{\GId_S}.\Complex \\
     \EvalM{\State}{\findmod{m}}{(\PMapDef, \Complex_S)}{F}}
    {\IntJ{\State}{\findmod{m}}{\kwlet\;n \mathbin{::} \Addr = \PMapDef}%
          {\State[\modules(m).\mathit{maps}(n) \mapsto
            (\mathsf{Ty}(\GId_S), F)]}}
\end{mathpar}

The address resolves to the source type $\GId_S$. The map definition is
evaluated with $\Complex_S$ as the source complex. For $\kwtotal$ maps, an
additional post-hoc check verifies that every generator in $\Complex_S$
has an image; the map is returned even if the check fails (errors
accumulate).

As with \textsc{TyLet}, name uniqueness is not checked at the type level.

% ---------------------------------------------------------------------------
\subsection{Module Inclusion}
\label{sec:type:import}
% ---------------------------------------------------------------------------

\begin{mathpar}
  \inferrule*[lab=\textsc{TyInclude}]
    {\mathit{alias} = n'\text{ if given, else }n \\
     \mathit{alias} \notin \mathit{used}(\findmod{m}) \\
     m' = \mathsf{resolve}(m, n) \\
     \IntJ{\State}{\cdot}{\Prog(m')}{\State_1} \\\\
     \Complex_{m'} = \State_1.\modules(m') \\
     F_{\mathit{id}} = \idmap{\Complex_{m'}} \\
     \mathit{gs} = \{ (\mathit{alias}.n_g, t_g, D_g) \mid
       (n_g, t_g, D_g) \in \Complex_{m'}.\mathit{gens},\;
       n_g \neq \varepsilon \} \\\\
     \State' = \State_1[
       \modules(m).\mathit{gens} \mathbin{{+}{=}} \mathit{gs},\;
       \modules(m).\mathit{maps}(\mathit{alias}) \mapsto
         (\mathsf{Mod}(m'), F_{\mathit{id}})]}
    {\IntJ{\State}{\findmod{m}}{\kwinclude\;n\;[\kwas\;n']}{\State'}}
\end{mathpar}

Module inclusion resolves $n$ to a canonical path $m'$, recursively
interprets $m'$ (which may transitively include further modules), and then
merges the included module's generators into the current module's complex.
The alias defaults to $n$ when no $\kwas$ clause is given; in either case,
the alias is checked for uniqueness in the current module.
Generators are prefixed with the alias and
duplicates---identified by tag, not by name---are silently skipped. This
is how diamond dependencies are handled: if two imports bring in the same
cell (same $\GId$), only the first copy is kept.

The state $\State_1$ from the recursive interpretation is carried forward
as the base for further modifications, ensuring that all types, cells, and
modules registered during the recursive call are visible.

% ===========================================================================
\section{Complex Resolution}
\label{sec:complex}
% ===========================================================================

Complex references resolve to namespaces $(\GId, \Complex)$:

\begin{mathpar}
  \inferrule*[lab=\textsc{Cpx-Root}]
    {\findmod{m}.\mathit{gens}(\varepsilon) = (\TagG{\GId_r}, \_, \_) \\
     \Complex = \findtype{\GId_r}.\Complex}
    {\ResJ{\State}{\varepsilon}{(\GId_r, \Complex)}}
  \and
  \inferrule*[lab=\textsc{Cpx-Addr}]
    {\AddrJ{\State}{\Addr}{\GId} \\
     \Complex = \findtype{\GId}.\Complex}
    {\ResJ{\State}{\Addr}{(\GId, \Complex)}}
  \and
  \inferrule*[lab=\textsc{Cpx-Block}]
    {\ResJ{\State}{\Addr}{(\GId, \Complex_0)} \\
     \IntJ{\State}{\Complex_0}{\overline{\CInst}}{(\State', \Complex')}}
    {\ResJ{\State}{\Addr\;\{\overline{\CInst}\}}{(\GId, \Complex')}}
\end{mathpar}

\textsc{Cpx-Root} handles the empty address by looking up the anonymous root
generator in the module complex. \textsc{Cpx-Addr} resolves a non-empty
address to a type's complex. \textsc{Cpx-Block} first resolves the
address, then interprets the body instructions, threading both state and
location; the body may add generators, diagrams, maps, includes, and
attaches. State changes from the body persist to the outer context.

\subsection{Address Resolution}
\label{sec:complex:addr}

\begin{mathpar}
  \inferrule*[lab=\textsc{Addr-Leaf}]
    {\Complex.\mathit{diags}(n) = D \\
     D\;\text{is a cell} \\
     \mathit{top}(D) = \TagG{\GId}}
    {\AddrJ{\State}{n}{\GId}}
  \and
  \inferrule*[lab=\textsc{Addr-Step}]
    {\Complex.\mathit{maps}(n_1) = (\mathsf{Mod}(m'), \_) \\
     \Complex' = \findmod{m'} \\
     \AddrJ{\State[\Complex']}
       {n_2.\cdots.n_k}{\GId}}
    {\AddrJ{\State}{n_1.n_2.\cdots.n_k}{\GId}}
\end{mathpar}

Address resolution navigates the \emph{map} namespace, not the generator
namespace. All segments except the last must be names of maps with module
domains; each step follows the module reference to its complex. The final
segment is looked up in the diagram namespace and must be a single cell.

\paragraph{Why maps, not generators?}
The map namespace is used because module inclusions register an identity
map under the alias name. This map's domain records which module (or type)
the inclusion came from, and its source complex provides the namespace for
further resolution.  Using maps rather than generators allows the same
resolution mechanism to work for both module paths
(\texttt{Theory.Equation}) and type paths (\texttt{F.S}).

% ===========================================================================
\section{Complex Body Instructions}
\label{sec:cinstr}
% ===========================================================================

Complex body instructions are interpreted in a mode $\mu \in \{\ModeG, \ModeL\}$.
In $\ModeG$ mode, new generators receive global tags and are registered in
$\State.\cells$. In $\ModeL$ mode, they receive local tags and are stored
in the complex's local cell store.

\subsection{Cell Generator}
\label{sec:cinstr:gen}

\begin{mathpar}
  \inferrule*[lab=\textsc{CGen}]
    {n \notin \mathit{used}(\Complex) \\
     \mathit{cd} = \text{(boundary from context)} \\
     \dim = \text{(inferred from } \mathit{cd}\text{)} \\\\
     \mu = \ModeG \implies
       \GId = \fresh,\;
       t = \TagG{\GId} \\
     \mu = \ModeL \implies
       t = \TagL{n} \\\\
     D = \CellOf{t}{\mathit{cd}} \\
     \Complex' = \Complex[\mathit{gens}(n) \mapsto D,\;
       \mathit{diags}(n) \mapsto D]}
    {\IntJ{\State}{\Complex}{n : \cdots}{\State', \Complex'}}
\end{mathpar}

A bare name declares a 0-cell ($\mathit{cd} = \CZero$, $\dim = 0$). A name
with boundary $n : d^{-} \to d^{+}$ declares an $n$-cell with
$\mathit{cd} = \CBdy{d^{-}}{d^{+}}$ and
$\dim = \mathsf{dim}(d^{-}) + 1$, where $d^{-}$ and $d^{+}$ are
evaluated as diagrams in the current complex.

The cell diagram $D = \CellOf{t}{\mathit{cd}}$ is constructed by taking
the pushout of the input and output boundary embeddings into a common lower
boundary, then adding one new top-dimensional cell. The resulting
ogposet (oriented graded poset; see \cref{sec:data})
has the combined shape of the two boundaries with the new cell on top.

The cell is entered into both $\mathit{gens}$ and $\mathit{diags}$.  In
$\ModeG$ mode the cell is also registered in $\State.\cells$; in $\ModeL$
mode it is added to $\Complex.\mathit{locals}$.

\subsection{Complex-Body Let and Def}
\label{sec:cinstr:let}

\begin{mathpar}
  \inferrule*[lab=\textsc{CLet}]
    {n \notin \mathit{used}(\Complex) \\
     \EvalD{\State}{\Complex}{\Diagram}{D}}
    {\IntJ{\State}{\Complex}{\kwlet\;n = \Diagram}
          {\State, \Complex[\mathit{diags}(n) \mapsto D]}}
  \and
  \inferrule*[lab=\textsc{CDef}]
    {n \notin \mathit{used}(\Complex) \\
     \AddrJ{\State}{\Addr}{\GId_S} \\
     \EvalM{\State}{\Complex}{(\PMapDef, \Complex_S)}{F}}
    {\IntJ{\State}{\Complex}{\kwlet\;n \mathbin{::} \Addr = \PMapDef}
          {\State, \Complex[\mathit{maps}(n) \mapsto (\mathsf{Ty}(\GId_S), F)]}}
\end{mathpar}

All \texttt{let} forms---type-level, complex-body, and local---check name
uniqueness.

\subsection{Include}
\label{sec:cinstr:include}

\begin{mathpar}
  \inferrule*[lab=\textsc{CInclude}]
    {\AddrJ{\State}{\Addr}{\GId} \\
     \Complex_S = \findtype{\GId}.\Complex \\
     n \notin \mathit{used}(\Complex) \\
     F_{\mathit{id}} = \idmap{\Complex_S} \\\\
     \mathit{gs} = \{ (n.n_g, t_g, D_g) \mid
       (n_g, t_g, D_g) \in \Complex_S.\mathit{gens},\;
       t_g \notin \Complex \} \\
     \Complex' = \Complex[\mathit{gens} \mathbin{{+}{=}} \mathit{gs},\;
       \mathit{maps}(n) \mapsto (\mathsf{Ty}(\GId), F_{\mathit{id}})]}
    {\IntJ{\State}{\Complex}{\kwinclude\;\Addr\;[\kwas\;n]}%
          {\State, \Complex'}}
\end{mathpar}

Include copies generators from the source type's complex into the current
complex, prefixing names, and adds an identity map. Crucially, no fresh
identifiers are allocated: the included generators retain their original
tags. This means that the ``same'' cell can appear in multiple complexes
under different names, but it has one canonical identity.

If no alias is given and the name would contain a dot (indicating a
non-local type), the interpreter raises an error: inclusion of non-local
types requires an explicit alias.

\subsection{Attach}
\label{sec:cinstr:attach}

\begin{mathpar}
  \inferrule*[lab=\textsc{CAttach}]
    {\AddrJ{\State}{\Addr}{\GId_A} \\
     \Complex_A = \findtype{\GId_A}.\Complex \\
     n \notin \mathit{used}(\Complex) \\
     F_0 = \text{(evaluate along map, or } \emptyset \text{ if omitted)} \\\\
     \text{for each } (n_g, t_g, \dim_g) \in \Complex_A.\mathit{gens}
       \text{, by ascending dimension, with } F_i(t_g) = \bot\text{:} \\
     \quad \GId' = \fresh \qquad
     \mathit{cd}' = \mathsf{bdy}(F_i, t_g) \qquad
     D' = \CellOf{\TagG{\GId'}}{\mathit{cd}'} \\
     \quad F_{i+1} = F_i[t_g \mapsto D'] \qquad
     \State_{i+1} = \State_i[\cells(\GId') \mapsto (\dim_g, \mathit{cd}')] \\\\
     \Complex' = \Complex[\mathit{gens} \mathbin{{+}{=}} \mathit{new},\;
       \mathit{maps}(n) \mapsto (\mathsf{Ty}(\GId_A), F_{\mathit{final}})]}
    {\IntJ{\State}{\Complex}{\kwattach\;n \mathbin{::} \Addr\;
       [\kwalong\;\PMapDef]}{\State', \Complex'}}
\end{mathpar}

\noindent
Here $\mathsf{bdy}(F, t)$ is $\CZero$ if the cell is 0-dimensional, and
$\CBdy{\apply{F}{d^{-}}}{\apply{F}{d^{+}}}$ otherwise (where
$\CBdy{d^{-}}{d^{+}}$ is the cell data of $t$).

Attachment differs from inclusion in a fundamental way: it creates
\emph{fresh} cells for every generator not already in the ``along'' map.
The ``along'' map specifies which generators of the attached type should
be identified with existing generators in the current complex (e.g.,
sharing a common base type), while every other generator gets a fresh copy.
Generators are processed in ascending dimension order, ensuring that when
we encounter an $n$-cell, all of its boundary cells have already been
mapped (either by the ``along'' map or by a previous iteration step).

\paragraph{Attach vs.\ include.}
Include reuses existing tags and creates no new cells; the result is an
aliased view of the same generators. Attach creates genuinely new cells
with fresh $\GId$s and new tag identities. This is the mechanism for
instantiating a ``template'': attaching the same type twice produces two
independent copies.

% ===========================================================================
\section{Local Instructions}
\label{sec:local}
% ===========================================================================

Local instructions operate within a resolved namespace
$(\GId_r, \Complex)$. Their defining feature is that definitions are
\emph{persisted}: the interpreter writes them back into the type's complex
in the global state.

\subsection{Local Let and Def}
\label{sec:local:let}

\begin{mathpar}
  \inferrule*[lab=\textsc{LocLet}]
    {n \notin \mathit{used}(\Complex) \\
     \EvalD{\State}{\Complex}{\Diagram}{D} \\
     D\;\text{has no local tags}}
    {\IntJ{\State}{(\GId_r, \Complex)}{\kwlet\;n = \Diagram}%
          {\State[\types(\GId_r).\Complex.\mathit{diags}(n) \mapsto D]}}
  \and
  \inferrule*[lab=\textsc{LocDef}]
    {n \notin \mathit{used}(\Complex) \\
     \EvalM{\State}{\Complex}{(\PMapDef, \Complex_S)}{F} \\
     F\;\text{has no local tags}}
    {\IntJ{\State}{(\GId_r, \Complex)}
      {\kwlet\;n \mathbin{::} \Addr = \PMapDef}%
      {\State[\types(\GId_r).\Complex.\mathit{maps}(n) \mapsto
        (\mathsf{Ty}(\GId_S), F)]}}
\end{mathpar}

Before persisting, the interpreter checks that the diagram or map
contains no locally-tagged cells. This ensures that definitions exported to
global state are self-contained. (This check is \emph{not} performed for
definitions inside complex bodies, which may legitimately refer to local
cells.)

The ``local'' in ``local block'' refers to the \emph{scope of operation}
(which type's complex we are working in), not to the lifetime of the
definition. Local-block definitions are global in duration.

\subsection{Assertions}
\label{sec:local:assert}

\begin{mathpar}
  \inferrule*[lab=\textsc{Assert-Diag}]
    {\EvalD{\State}{\Complex}{\Diagram_1}{D_1} \\
     \EvalD{\State}{\Complex}{\Diagram_2}{D_2} \\
     D_1 \cong D_2}
    {\AssertJ{\State}{\Complex}{\kwassert\;\Diagram_1 = \Diagram_2}}
  \and
  \inferrule*[lab=\textsc{Assert-Map}]
    {\EvalT{\State}{\Complex}{\Diagram_1}{\MTerm{F_1, \Complex_S}} \\
     \EvalT{\State}{\Complex}{\Diagram_2}{\MTerm{F_2, \Complex_S}} \\\\
     \forall (n, t) \in \Complex_S.\mathit{gens}{:}\;
       \mathsf{def}(F_1, t) = \mathsf{def}(F_2, t) \\
     \forall t\;\text{defined in both}{:}\;
       F_1(t) \cong F_2(t)}
    {\AssertJ{\State}{\Complex}{\kwassert\;\Diagram_1 = \Diagram_2}}
\end{mathpar}

Both sides must evaluate to terms of the same kind (both diagram or both
map). Diagram equality is checked via isomorphism of the underlying
ogposets with matching labels. Map equality is checked pointwise:  for
each generator in the shared source complex, both maps must agree on
definedness and have isomorphic images.

\paragraph{Structural, not propositional.}
Assertions check \emph{definitional} (structural) equality---isomorphism
of the underlying ogposets with matching labels. There are no non-trivial
equalities: everything interesting is a \emph{rewrite} (a
higher-dimensional cell). For example, the unit law
$\kwassert\;\mathit{id}\;{}\alpha = \alpha$ will \emph{fail}, because
$\mathit{id}\;{}\alpha$ has an additional identity cell in its ogposet
that $\alpha$ alone does not. To express the unit law, one declares a
cell $\lambda : \mathit{id}\;{}\alpha \to \alpha$, which is itself a
rewrite witnessing the relationship.

However, middle-four interchange---the equation
$(\alpha \paste{0} \beta)(\gamma \paste{0} \delta) =
(\alpha\gamma) \paste{0} (\beta\delta)$---\emph{is} structural: both
sides produce isomorphic ogposets.

% ===========================================================================
\section{Diagram Evaluation}
\label{sec:diag}
% ===========================================================================

\subsection{Implicit Pasting}
\label{sec:diag:juxt}

\begin{mathpar}
  \inferrule*[lab=\textsc{Diag-One}]
    {\EvalT{\State}{\Complex}{\DExpr}{\DTerm{D}}}
    {\EvalD{\State}{\Complex}{\DExpr}{D}}
  \and
  \inferrule*[lab=\textsc{Diag-Juxt}]
    {\EvalD{\State}{\Complex}{\DExpr_1\;\cdots\;\DExpr_{k-1}}{D_L} \\
     \EvalD{\State}{\Complex}{\DExpr_k}{D_R} \\\\
     j = \min(\dim(D_L), \dim(D_R)) - 1 \\
     D = D_L \paste{j} D_R}
    {\EvalD{\State}{\Complex}{\DExpr_1\;\cdots\;\DExpr_k}{D}}
\end{mathpar}

Juxtaposition is left-associative. The Rust implementation evaluates
implicit pasting left to right: the first expression is evaluated, then
each subsequent expression is evaluated and pasted onto the accumulator.
This differs from explicit pasting, which evaluates the right side before
the left (\cref{sec:diag:paste}). The paste dimension is
$\min(\dim(D_L), \dim(D_R)) - 1$. For two 1-cells this gives $j = 0$
(horizontal composition); for two 2-cells, $j = 1$ (vertical
composition).

\subsection{Explicit Pasting}
\label{sec:diag:paste}

\begin{mathpar}
  \inferrule*[lab=\textsc{Diag-Paste}]
    {\EvalD{\State}{\Complex}{\overline{\DExpr}_R}{D_R} \\
     \EvalD{\State}{\Complex}{\Diagram_L}{D_L} \\
     \bdy{+}{k}{D_L} \cong \bdy{-}{k}{D_R} \\
     D = D_L \paste{k} D_R}
    {\EvalD{\State}{\Complex}{\Diagram_L \paste{k}
      \overline{\DExpr}_R}{D}}
\end{mathpar}

The right-hand side is evaluated before the left-hand side, matching
the OCaml reference implementation's evaluation order.
The paste is defined only when $D_L$'s output $k$-boundary
matches $D_R$'s input $k$-boundary. The result is computed as a pushout
of the two boundary embeddings, merging labels and paste trees.

\subsection{Component Evaluation}
\label{sec:diag:comp}

\begin{mathpar}
  \inferrule*[lab=\textsc{Comp-Name-D}]
    {\Complex.\mathit{diags}(n) = D}
    {\EvalT{\State}{\Complex}{n}{\DTerm{D}}}
  \and
  \inferrule*[lab=\textsc{Comp-Name-M}]
    {\Complex.\mathit{diags}(n) = \bot \\
     \Complex.\mathit{maps}(n) = (\mathit{dom}, F) \\
     \Complex_S = \mathsf{src}(\State, \mathit{dom})}
    {\EvalT{\State}{\Complex}{n}{\MTerm{F, \Complex_S}}}
\end{mathpar}

\begin{mathpar}
  \inferrule*[lab=\textsc{Comp-Paren}]
    {\EvalD{\State}{\Complex}{\Diagram}{D}}
    {\EvalT{\State}{\Complex}{(\Diagram)}{\DTerm{D}}}
  \and
  \inferrule*[lab=\textsc{Comp-PMap}]
    {\EvalM{\State}{\Complex}{(\PMapG, \Complex)}{F, \Complex_S}}
    {\EvalT{\State}{\Complex}{(\PMapG)}{\MTerm{F, \Complex_S}}}
\end{mathpar}

\begin{mathpar}
  \inferrule*[lab=\textsc{Comp-AnonMap}]
    {\ResJ{\State}{\Cpx}{(\_, \Complex_T)} \\
     \EvalM{\State}{\Complex_T}{(\PMapDef, \Complex)}{F}}
    {\EvalT{\State}{\Complex}{(\,\kwmap\;\PMapDef :: \Cpx\,)}
      {\MTerm{F, \Complex}}}
  \and
  \inferrule*[lab=\textsc{Comp-Bd}]
    {\strut}
    {\EvalT{\State}{\Complex}{\kwin}{\mathsf{Bd}(\text{in})}}
\end{mathpar}

Name lookup checks diagrams first, then maps.

\paragraph{Anonymous maps in diagram context.}
Anonymous maps resolve $\Cpx$ to obtain $\Complex_T$, which
becomes the evaluation context for clause right-hand sides (images).
Clause left-hand sides are resolved in the \emph{ambient} complex
$\Complex$, which also becomes the returned map's source. In other words,
the $\mathbin{::}\Cpx$ annotation in a diagram-level anonymous map
specifies where images live, not where sources come from. This differs
from the convention in named map definitions ($\kwlet\;n \mathbin{::}
\Addr = \ldots$), where $\Addr$ identifies the source type.

The source complex for
$\mathsf{src}(\State, \mathit{dom})$ is
$\findtype{\GId}.\Complex$ if $\mathit{dom} = \mathsf{Ty}(\GId)$, or
$\findmod{m}$ if $\mathit{dom} = \mathsf{Mod}(m)$.

\subsection{Dot Access}
\label{sec:diag:dot}

The dot operator has three semantically distinct behaviours:

\begin{mathpar}
  \inferrule*[lab=\textsc{Dot-Bd}]
    {\EvalT{\State}{\Complex}{\DExpr}{\DTerm{D}} \\
     k = \max(\dim(D) - 1, 0) \\
     D' = \bdy{s}{k}{D}}
    {\EvalT{\State}{\Complex}{\DExpr.\mathsf{Bd}(s)}{\DTerm{D'}}}
\end{mathpar}

\begin{mathpar}
  \inferrule*[lab=\textsc{Dot-Apply}]
    {\EvalT{\State}{\Complex}{\DExpr}{\MTerm{F, \Complex_S}} \\
     \EvalT{\State}{\Complex_S}{\DComp}{\DTerm{D}} \\
     D' = \apply{F}{D}}
    {\EvalT{\State}{\Complex}{\DExpr.\DComp}{\DTerm{D'}}}
\end{mathpar}

\begin{mathpar}
  \inferrule*[lab=\textsc{Dot-Compose}]
    {\EvalT{\State}{\Complex}{\DExpr}{\MTerm{F, \Complex_S}} \\
     \EvalT{\State}{\Complex_S}{\DComp}{\MTerm{G, \Complex_{S'}}} \\
     H = \compose{F}{G}}
    {\EvalT{\State}{\Complex}{\DExpr.\DComp}
      {\MTerm{H, \Complex_{S'}}}}
\end{mathpar}

\textsc{Dot-Bd}: extracting the codimension-1 boundary of a diagram.

\textsc{Dot-Apply}: the key rule. The right-hand component is evaluated
in the source complex $\Complex_S$ of the map, not in the ambient
complex $\Complex$. This \emph{context switch} is how \texttt{F.x}
resolves \texttt{x} in $F$'s domain rather than in the surrounding
namespace. The result is the image of $D$ under $F$.

\textsc{Dot-Compose}: when both sides are maps, compose them. The
source complex of the result is $\Complex_{S'}$ (the inner map's source).

\paragraph{Dot on diagram with a non-boundary component} produces an error.
Similarly, a boundary accessor on a map term produces an error. These
cases are intentionally not given rules: they signal a type error in the
user's program.

% ===========================================================================
\section{Partial Map Evaluation}
\label{sec:pmap}
% ===========================================================================

Map expressions evaluate relative to a source complex $\Complex_S$
(the domain of the map being defined). We write the judgment as
$\EvalM{\State}{\Complex}{(\PMapDef, \Complex_S)}{F}$.

\subsection{Basic Lookup and Composition}
\label{sec:pmap:basic}

\begin{mathpar}
  \inferrule*[lab=\textsc{PM-Name}]
    {\Complex.\mathit{maps}(n) = (\mathit{dom}, F) \\
     \Complex_S = \mathsf{src}(\State, \mathit{dom})}
    {\EvalM{\State}{\Complex}{n}{\MC{F}{\Complex_S}}}
  \and
  \inferrule*[lab=\textsc{PM-Dot}]
    {\EvalM{\State}{\Complex}{\PMapBasic}{\MC{F}{\Complex_S}} \\
     \EvalM{\State}{\Complex_S}{\PMapG}{\MC{G}{\Complex_{S'}}} \\
     H = \compose{F}{G}}
    {\EvalM{\State}{\Complex}{\PMapBasic.\PMapG}{\MC{H}{\Complex_{S'}}}}
\end{mathpar}

Map composition via dot follows the same pattern as diagram dot access:
the right side is resolved in the left side's source complex.

\subsection{Extension Maps}
\label{sec:pmap:ext}

\begin{mathpar}
  \inferrule*[lab=\textsc{PM-Ext}]
    {F_0 = \text{(prefix map, or } \emptyset \text{)} \\
     \text{for each clause } \Diagram_i \Rightarrow \Diagram'_i\text{:} \\
     \quad \EvalD{\State}{\Complex_S}{\Diagram_i}{D_i} \qquad
     \EvalD{\State}{\Complex}{\Diagram'_i}{D'_i} \\
     \quad F_{i+1} = \smartext(F_i, \Complex_S, \Complex, D_i, D'_i)}
    {\EvalM{\State}{\Complex}{(\PMapExt, \Complex_S)}{F_{\mathit{final}}}}
\end{mathpar}

In each clause, the left-hand side is evaluated in the source complex
$\Complex_S$ and the right-hand side in the current location $\Complex$.
This is because each clause maps a cell of the source type to a diagram
in the ambient complex. Clauses are processed sequentially: each extends
the map built by the previous, so later clauses see earlier mappings.

\subsection{The Smart Extension Algorithm}
\label{sec:pmap:smart}

The function $\smartext(F, \Complex_S, \Complex_T, D_s, D_t)$ extends
a partial map $F$ with a mapping from a source cell $D_s$ to a target
diagram $D_t$, automatically inferring mappings for unmapped boundary
cells:

\begin{enumerate}
  \item \textbf{Cell check.} $D_s$ must be a single cell. Extract its
    tag $t$.
  \item \textbf{Idempotence.} If $F(t)$ is already defined and
    isomorphic to $D_t$, return $F$ unchanged. If defined but not
    isomorphic, raise an error.
  \item \textbf{Boundary collection.} Look up the cell data for $t$
    (checking $\State.\cells$ first, then $\State.\types$ as fallback,
    since types also carry cell data). If
    it has boundaries $\CBdy{d^{-}}{d^{+}}$, collect all top-dimensional
    tags in both that are not yet mapped by $F$.
  \item \textbf{Recursive extension.} For each unmapped boundary tag $t_b$
    with sign $s$:
    \begin{enumerate}
      \item Compute $D_b = \bdy{s}{\dim-1}{D_t}$ (the corresponding
        boundary of the \emph{target}).
      \item If the source boundary is a single cell, recurse.
      \item If the source boundary is composite, find an isomorphism
        between the source and target boundary shapes, use it to
        identify which target tag corresponds to $t_b$, then recurse.
    \end{enumerate}
  \item \textbf{Validation.} Call $\mathsf{PMap{:}{:}extend}$ to add
    $t \mapsto D_t$, which checks boundary compatibility.
\end{enumerate}

\paragraph{Isomorphism-based inference.}
When the source boundary is composite, the algorithm must determine which
cell in the target boundary corresponds to each unmapped source cell.
It normalises both boundary ogposets, finds a bijection between their
cells, and uses this to transfer labels. If a source tag appears at
multiple positions in the boundary, the algorithm checks that all
corresponding target positions carry the same label. This is the mechanism
that makes terse definitions like \texttt{[merge => split]} automatically
infer boundary mappings: the user maps only the top-dimensional cells,
and the algorithm deduces where the boundary cells must go by matching
boundary shapes.

\subsection{Map-to-Map Assignment}
\label{sec:pmap:mapassign}

When a clause maps a map term to another map term (both $\MTerm{\cdot}$),
the interpreter processes it as a pointwise extension:
\begin{enumerate}
  \item Verify that both maps share the same source complex.
  \item Iterate over generators of the shared source, sorted by dimension.
  \item For each generator defined in both maps: if the left image is a
    cell, call $\smartext$ to map the left image to the right image. If
    the left image is composite, check that all its labels are already
    defined in the current map (but do not recurse).
  \item A generator defined in one map but not the other is an error.
\end{enumerate}

This handles constructions like \texttt{[Dom => Ob, Cod => Ob]} where
\texttt{Dom} and \texttt{Cod} resolve to maps.

% ===========================================================================
\section{Map Application and Composition}
\label{sec:apply}
% ===========================================================================

The operation $\apply{F}{D}$ applies a partial map to a diagram. There
are two paths:

\begin{itemize}
  \item \textbf{Cellular.} If $F$ is cellular (all images are single
    cells), simply remap each label tag in $D$ to the corresponding image
    tag. The ogposet shape is shared. This is the common case for identity
    maps and inclusion maps.
  \item \textbf{General.} Follow the paste tree of $D$ recursively. For a
    $\mathsf{Leaf}(t)$, return $F(t)$. For a $\mathsf{Node}(k, L, R)$,
    recursively apply to $L$ and $R$, then paste the results at dimension
    $k$.
\end{itemize}

The identity map $\idmap{\Complex}$ maps every generator's tag to its own
classifier diagram. It is always cellular and therefore enjoys the fast
application path.

Map composition $\compose{F}{G}$ is computed eagerly: for each entry
$t \mapsto D$ in $G$'s table, the result maps $t$ to $\apply{F}{D}$.

% ===========================================================================
\section{Data Structures}
\label{sec:data}
% ===========================================================================

For completeness, we briefly describe the key structures underlying
diagram manipulation.

\subsection{Oriented Graded Posets}
An oriented graded poset (ogposet) $G$ has a dimension, a set of cells at
each dimension, and face maps $\partial^{\pm}_d : G_d \to \mathcal{P}(G_{d-1})$ giving the input and output faces. Key operations:
normalisation (canonical cell ordering, cached by pointer identity),
boundary extraction ($\bdy{\pm}{k}{G}$, also cached), and isomorphism
finding (normalise both, compose embeddings).

\subsection{Diagrams}
A diagram $D$ is an ogposet (the shape) together with labels
$\ell_d : G_d \to \textsf{Tag}$ and paste trees recording how the
diagram was built. The distinction between a cell and a composite is
recorded in the paste tree: a cell has a $\mathsf{Leaf}$ at the top level.

\subsection{Partial Maps}
A partial map $F$ is a table $\textsf{Tag} \rightharpoonup
(\CData, \Diag{D})$, with a Boolean \texttt{cellular} flag. Tags are
indexed by dimension for ordered iteration. The $\mathsf{extend}$ operation
validates boundary compatibility before adding a new entry.

% ===========================================================================
\section{Summary of Subtleties}
\label{sec:pain}
% ===========================================================================

\begin{enumerate}
  \item \textbf{Name ambiguity.} A name may resolve to a diagram or a
    map. Diagrams are checked first. The dot operator dispatches on the
    left operand's kind.

  \item \textbf{Evaluation order.} Implicit pasting is left-associative
    and evaluates left to right.
    Explicit pasting ($\paste{k}$) evaluates the right side first,
    matching the OCaml reference implementation.

  \item \textbf{Source complex switch.} Dotting into a map evaluates the
    right component in the map's source complex, not the ambient complex.
    This is the mechanism behind \texttt{F.x}.

  \item \textbf{Local block persistence.} Definitions in local blocks are
    persisted to global state. ``Local'' refers to scope of operation, not
    lifetime.

  \item \textbf{No local tags in persisted definitions.} Local blocks
    check for local tags before persisting. Complex bodies do not.

  \item \textbf{Attach creates fresh cells; include reuses tags.} Attach
    processes generators by ascending dimension to ensure boundary
    computability.

  \item \textbf{Smart extension.} The $\smartext$ algorithm infers
    boundary mappings via ogposet isomorphisms. This can fail if boundary
    shapes differ or if a tag maps inconsistently.

  \item \textbf{Totality is post-hoc.} Checked after all clauses,
    non-fatal.

  \item \textbf{Module idempotency.} A module is interpreted at most once.
    Import order can affect visible state.

  \item \textbf{Structural equality only.} Unit laws fail; middle-four
    interchange passes. Non-trivial equalities are higher cells.

  \item \textbf{Self-referencing identity.} A type's name inside its own
    complex resolves to the identity inclusion map.

  \item \textbf{Diamond imports.} Duplicate generators (same tag) from
    multiple imports are silently deduplicated.

  \item \textbf{Anonymous map convention.} In a named map definition
    $\kwlet\;f \mathbin{::} T = \ldots$, the annotation $T$ identifies
    the \emph{source}. In a diagram-level anonymous map
    $(\kwmap\;\ldots \mathbin{::} T)$, the annotation $T$ identifies
    where \emph{images} are evaluated---the opposite convention.

  \item \textbf{Type-level name uniqueness.} Name collisions are checked
    in complex bodies and local blocks but not in type blocks, where a
    duplicate $\kwlet$ silently overwrites.
\end{enumerate}

\end{document}
